% =============================================================================
%  DOCUMENTACAO DO TRABALHO DE REDES DE COMPUTADORES - UEMS
%  Plataforma de Monitoramento Inteligente de Filas em Tempo Real
%  Aluno: Victor Vendruscolo
%
%  COMO COMPILAR
%    Motor: pdfLaTeX. No Overleaf, e o padrao (Menu > Compiler > pdfLaTeX).
%    Compile duas vezes para o sumario e as referencias cruzadas fecharem.
%
%  IMAGENS
%    Coloque as capturas de tela na subpasta figuras/, com os nomes exatos
%    listados em figuras/LEIA-ME.txt. Enquanto um arquivo nao existir, o
%    documento compila e mostra uma moldura no lugar, com o nome esperado.
% =============================================================================

\documentclass[12pt,a4paper]{article}

% --- Idioma e codificacao ---------------------------------------------------
\usepackage[utf8]{inputenc}
\usepackage[T1]{fontenc}
\usepackage[brazil]{babel}

% --- Fonte: Times, para acompanhar o documento do enunciado -----------------
\usepackage{mathptmx}

% --- Layout -----------------------------------------------------------------
\usepackage[top=2.5cm,bottom=2.5cm,left=3cm,right=2.5cm]{geometry}
\usepackage{indentfirst}
\setlength{\parindent}{1.25cm}
\setlength{\parskip}{3pt}
% Permite quebra de linha dentro de nomes longos em fonte de maquina de
% escrever (\texttt), evitando linhas que estouram a margem.
\usepackage[htt]{hyphenat}
\sloppy

% --- Tabelas, figuras e codigo ---------------------------------------------
\usepackage{booktabs}
\usepackage{array}
\usepackage{longtable}
\usepackage{graphicx}
\usepackage{float}
\usepackage{listings}
\usepackage[font=small,labelfont=bf]{caption}
\usepackage{ragged2e}
\usepackage{url}
\usepackage[hidelinks]{hyperref}

% Coluna justificada de largura fixa, usada nas tabelas longas.
\newcolumntype{L}[1]{>{\RaggedRight\arraybackslash}p{#1}}

% --- Estilo dos trechos de codigo ------------------------------------------
\lstset{
  basicstyle=\ttfamily\footnotesize,
  keywordstyle=\bfseries,
  commentstyle=\itshape,
  language=C,
  numbers=none,
  frame=single,
  framesep=6pt,
  breaklines=true,
  showstringspaces=false,
  captionpos=b,
  xleftmargin=0.5cm,
  xrightmargin=0.5cm
}

% --- Macro das capturas de tela ---------------------------------------------
% Se o arquivo existir em figuras/, ele e inserido. Se nao existir, aparece
% uma moldura com o nome esperado, e o documento continua compilando.
\newcommand{\tela}[3]{%
  \begin{figure}[H]
    \centering
    \IfFileExists{figuras/#1}{%
      \includegraphics[width=#3\textwidth]{figuras/#1}%
    }{%
      \fbox{%
        \begin{minipage}[c][4.5cm][c]{0.85\textwidth}
          \centering\ttfamily\small
          [CAPTURA PENDENTE]\\[6pt]
          figuras/#1
        \end{minipage}%
      }%
    }
    \caption{#2}
  \end{figure}%
}

\begin{document}

% =============================================================================
% CAPA - segue o layout do enunciado (trabalho_redes_2026.pdf), com o
% acrescimo do nome do aluno, que o enunciado nao traz por ser o documento
% do professor.
% =============================================================================
\begin{titlepage}
\thispagestyle{empty}
\begin{center}

{\large Universidade Estadual de Mato Grosso do Sul}\\[10pt]
{\large Curso de Ciência da Computação}\\[10pt]
{\large Disciplina de Redes de Computadores}\\[10pt]
{\large Trabalho de Redes de Computadores}

\vspace{7cm}

{\large \underline{\parbox{\textwidth}{\centering PLATAFORMA DE MONITORAMENTO INTELIGENTE DE FILAS\\ EM TEMPO REAL}}}

\vspace{2cm}

{\large Documentação}

\vspace{1.5cm}

{\large Victor Vendruscolo}

\vfill

{\large Dourados (MS), 29 de julho de 2026.}

\end{center}
\end{titlepage}

% =============================================================================
% SUMARIO
% =============================================================================
\thispagestyle{empty}
\tableofcontents
\newpage

\setcounter{page}{1}

% =============================================================================
% ITEM 1 - SUMARIO DO PROBLEMA
% =============================================================================
\section{Sumário do problema}

Este trabalho implementa uma plataforma cliente-servidor para monitorar uma fila
de atendimento em tempo real. Um servidor central mantém a fila; vários clientes
se conectam ao mesmo tempo e operam sobre ela. O sistema foi escrito em C, sobre
sockets TCP, e roda em Linux.

O problema tratado é a falta de informação compartilhada em locais de
atendimento ao público. Quando cada posto anota a fila por conta própria,
ninguém sabe o estado real do atendimento. O enunciado cita superlotação, filas
desorganizadas e ausência de atualização em tempo real como sintomas disso.

A solução adotada centraliza a fila em um único servidor. Todo cliente que
cadastra um usuário altera a mesma estrutura, e todos os demais são avisados na
hora. Assim, o estado da fila é único e visível para todos os operadores
conectados.

O sistema oferece quatro operações, que reproduzem as telas do enunciado:

\begin{itemize}
  \item \textbf{Autenticação}: o cliente se identifica ao conectar, com uma
        credencial fixa definida no código.
  \item \textbf{Adicionar usuário}: o operador informa um identificador e um
        nome, que entram na fila compartilhada.
  \item \textbf{Ver fila}: o servidor devolve a fila completa, do primeiro ao
        último registro.
  \item \textbf{\textit{Heartbeat}}: o servidor devolve apenas os usuários
        cadastrados por outros clientes desde a última verificação. Se não
        houver novidade, devolve \texttt{ALIVE}.
\end{itemize}

Além dessas quatro operações pedidas, o servidor envia uma notificação
espontânea. Sempre que alguém cadastra um usuário, todos os outros clientes
recebem a mensagem de \textit{broadcast} sem terem pedido nada. É esse
mecanismo que caracteriza o tempo real do sistema.

\vspace{6pt}
\noindent\textbf{Terminologia.} As telas do enunciado usam a palavra
``paciente'', mas o texto do enunciado usa ``usuário''. Este documento adota
\textbf{usuário} para quem entra na fila e \textbf{cliente} para o programa
\texttt{./cliente} que se conecta e opera o sistema. São dois conceitos
diferentes e a distinção vale para todo o documento.

% =============================================================================
% ITEM 2 - ALGORITMOS, TADs, FUNCOES E DECISOES
% =============================================================================
\section{Algoritmos, tipos de dados, funções e decisões}

\subsection{Arquitetura geral}

O sistema tem dois programas, ligados por uma conexão TCP na porta 8080. A
persistência fica em arquivos de texto, do lado do servidor.

\begin{figure}[H]
\centering
\begin{minipage}{0.9\textwidth}
\ttfamily\footnotesize
\begin{verbatim}
                       porta 8080 (TCP)
   [Cliente 1] ------------+
   [Cliente 2] ------------+---> [Servidor] --- fila compartilhada (mutex)
   [Cliente N] ------------+           |
                                       +---> dados/*.txt (persistencia)
\end{verbatim}
\end{minipage}
\caption{Arquitetura do sistema, correspondente ao diagrama do enunciado.}
\end{figure}

A técnica de múltiplos clientes escolhida foi \textbf{multithreads}
(\texttt{pthreads}), uma das três permitidas pelo enunciado. O servidor cria uma
thread para cada conexão aceita. Duas razões pesaram nessa escolha. A primeira é
que a fila é compartilhada por todos os clientes: com threads, todas enxergam a
mesma memória, e basta um \textit{mutex} para proteger o acesso. A segunda é que
o \textit{broadcast} precisa alcançar sockets de outras conexões, o que também
sai de graça no mesmo processo. Com \texttt{fork()}, as duas coisas exigiriam
comunicação entre processos.

A thread principal do servidor não atende nenhum cliente. Ela apenas aceita
conexões, cria a thread correspondente e volta a aceitar. Isso atende
literalmente ao enunciado, que reserva ao servidor principal apenas aceitar
conexões, distribuir eventos e monitorar sockets.

\begin{lstlisting}[caption={Laço principal do servidor, em \texttt{servidor.c}. A thread criada recebe uma \texttt{struct} própria, com o socket, o IP de origem e o número da conexão.}]
for (;;) {
    int  sock_cliente = accept(sock_escuta, ...);
    ...
    ArgumentoThread *arg = malloc(sizeof(ArgumentoThread));
    arg->sock   = sock_cliente;
    arg->numero = ++total_conexoes;
    pthread_create(&thread, &atributos, atende_cliente, arg);
}
\end{lstlisting}

\subsection{Tipo abstrato de dados: a fila}

A fila é o dado central do sistema. Ela é um vetor de tamanho fixo, protegido
por um único \textit{mutex}, e todo acesso passa pelas funções do módulo
\texttt{fila.c}. Manter o acesso concentrado em um módulo torna impossível
esquecer de travar.

\begin{lstlisting}[caption={Tipo \texttt{Usuario} e estado interno da fila.}]
typedef struct {
    int  id;                    /* identificador digitado pelo operador */
    char nome[TAM_NOME];        /* nome, terminado em '\0'              */
} Usuario;

static Usuario         fila[MAX_USUARIOS_FILA];  /* 20000 registros */
static int             total;                    /* posicoes em uso */
static pthread_mutex_t mutex;                    /* protege os dois */
\end{lstlisting}

O vetor foi preferido à lista ligada pela simplicidade. O limite de 20.000
registros cobre com folga o teste de carga de 10.000 clientes, mesmo que cada um
insira mais de um usuário. Esse limite é uma decisão de implementação declarada
na Seção~\ref{sec:omissos}.

Cada conexão guarda, por conta própria, um marcador do último índice da fila que
já viu. É esse marcador que permite ao \textit{heartbeat} devolver apenas o que
é novo para aquele cliente.

\subsection{Protocolo de mensagens}
\label{sec:protocolo}

Toda mensagem é uma linha de texto ASCII terminada em \texttt{\textbackslash n}.
O comando vem em maiúsculas, seguido dos argumentos separados por espaço. Essa
escolha resolve o problema central de trabalhar sobre TCP: o TCP entrega um
fluxo de bytes, não mensagens delimitadas, e uma chamada de \texttt{send()} não
corresponde necessariamente a uma de \texttt{recv()}. Quem lê acumula os bytes em
um buffer próprio até encontrar a quebra de linha.

Respostas de várias linhas usam um cabeçalho e um rodapé fixos. O cliente lê
linha a linha até bater no rodapé, sem precisar saber quantas linhas virão. É o
formato que o próprio enunciado mostra nas telas.

\begin{table}[H]
\centering
\caption{Mensagens do cliente para o servidor.}
\footnotesize
\begin{tabular}{L{3.2cm} L{5.6cm} L{4.4cm}}
\toprule
\textbf{Operação} & \textbf{Formato} & \textbf{Quando} \\
\midrule
Login & \texttt{LOGIN <usuario> <senha>} & Primeira mensagem da sessão \\
Adicionar usuário & \texttt{ADD <seq> <id> <nome>} & Opção \texttt{1} do menu \\
Ver fila & \texttt{LIST <seq>} & Opção \texttt{2} do menu \\
\textit{Heartbeat} & \texttt{HEARTBEAT <seq>} & Opção \texttt{3} do menu \\
Sair & \texttt{SAIR} & Opção \texttt{0} do menu \\
\bottomrule
\end{tabular}
\end{table}

\begin{table}[H]
\centering
\caption{Mensagens do servidor para o cliente.}
\footnotesize
\begin{tabular}{L{3.2cm} L{5.6cm} L{4.4cm}}
\toprule
\textbf{Resposta} & \textbf{Formato} & \textbf{Quando} \\
\midrule
Login aceito & \texttt{LOGIN\_OK} & A credencial confere \\
Login recusado & \texttt{LOGIN\_FAIL} & A credencial não confere; o servidor fecha a conexão \\
Usuário inserido & \texttt{ADD\_OK <id> <nome>} & Após gravar na fila e disparar o \textit{broadcast} \\
Fila & \texttt{===== FILA =====} \newline \texttt{<id> - <nome>} $\times$ N \newline \texttt{================} & Resposta de \texttt{LIST} \\
\textit{Heartbeat} com novidade & Mesmo formato da fila, só com os registros de outros clientes & Resposta de \texttt{HEARTBEAT} \\
\textit{Heartbeat} sem novidade & \texttt{ALIVE} & Resposta de \texttt{HEARTBEAT} \\
Erro & \texttt{ERRO <descrição>} & Comando inválido; a conexão continua aberta \\
\midrule
\textit{Broadcast} & \texttt{[Broadcast] Novo usuario: <nome>} & Mensagem não solicitada, enviada a todos os clientes menos ao autor do \texttt{ADD} \\
\bottomrule
\end{tabular}
\end{table}

O número de sequência \texttt{<seq>} acompanha todo comando posterior ao login.
Ele começa em 1 e aumenta a cada novo comando daquela conexão. Sua função é
sustentar a retransmissão, tratada na Seção~\ref{sec:retransmissao}.

\subsection{Módulos e funções}

O código está dividido em seis módulos, além dos dois programas principais e do
gerador de carga. Cada arquivo tem uma responsabilidade só.

{\footnotesize
\begin{longtable}{L{2.8cm} L{4.1cm} L{6.3cm}}
\caption{Módulos do sistema e suas funções principais.}\\
\toprule
\textbf{Arquivo} & \textbf{Responsabilidade} & \textbf{Funções principais} \\
\midrule
\endfirsthead
\toprule
\textbf{Arquivo} & \textbf{Responsabilidade} & \textbf{Funções principais} \\
\midrule
\endhead
\texttt{comum.h} &
Constantes, tipo \texttt{Usuario} e textos do protocolo &
Somente definições \\
\addlinespace
\texttt{protocolo.c/.h} &
Envio e leitura de mensagens delimitadas por linha &
\texttt{protocolo\_leitor\_init}, \texttt{protocolo\_le\_linha}, \texttt{protocolo\_envia\_linha}, \texttt{protocolo\_limpa\_bordas} \\
\addlinespace
\texttt{fila.c/.h} &
Fila compartilhada protegida por \textit{mutex} &
\texttt{fila\_init}, \texttt{fila\_adiciona}, \texttt{fila\_tamanho}, \texttt{fila\_copia\_intervalo} \\
\addlinespace
\texttt{sessoes.c/.h} &
Registro dos clientes conectados e envio do \textit{broadcast} &
\texttt{sessoes\_init}, \texttt{sessoes\_registra}, \texttt{sessoes\_remove}, \texttt{sessoes\_broadcast}, \texttt{sessoes\_trava\_envio}, \texttt{sessoes\_libera\_envio} \\
\addlinespace
\texttt{persistencia.c/.h} &
Gravação nos quatro arquivos de texto &
\texttt{persistencia\_init}, \texttt{persistencia\_log\_servidor}, \texttt{persistencia\_log\_sessao}, \texttt{persistencia\_historico\_add}, \texttt{persistencia\_salva\_fila} \\
\addlinespace
\texttt{servidor.c} &
Aceitação de conexões e atendimento &
\texttt{main}, \texttt{atende\_cliente}, \texttt{autentica}, \texttt{trata\_add}, \texttt{trata\_list}, \texttt{trata\_heartbeat}, \texttt{envia\_bloco\_fila}, \texttt{envia\_resposta}, \texttt{verifica\_sequencia}, \texttt{guarda\_resposta\_add}, \texttt{nome\_valido} \\
\addlinespace
\texttt{cliente.c} &
Interface de operação &
\texttt{main}, \texttt{conecta\_ao\_servidor}, \texttt{autentica}, \texttt{thread\_receptora}, \texttt{envia\_comando}, \texttt{imprime\_resposta}, \texttt{opcao\_adicionar\_usuario}, \texttt{opcao\_ver\_fila}, \texttt{opcao\_heartbeat}, \texttt{mostra\_menu} \\
\addlinespace
\texttt{carga.c} &
Gerador automático de clientes para o teste de carga &
\texttt{main}, \texttt{conecta}, \texttt{autentica} \\
\addlinespace
\texttt{Makefile} &
Compilação &
Alvos \texttt{all} (padrão), \texttt{carga}, \texttt{tudo}, \texttt{clean} \\
\bottomrule
\end{longtable}}

O \texttt{Makefile} sem parâmetros gera exatamente os dois programas exigidos,
\texttt{cliente} e \texttt{servidor}. O gerador de carga ficou em um alvo
separado, \texttt{make carga}, justamente para não acrescentar um terceiro
executável ao alvo padrão.

Não existem \texttt{servidor.h} nem \texttt{cliente.h}. Um cabeçalho não faz
sentido para um módulo que só contém o \texttt{main} e funções privadas ao
arquivo, todas declaradas \texttt{static}.

\subsection{Algoritmos principais}

\subsubsection*{Atendimento de um cliente}

Cada thread executa sempre a mesma sequência. Ela autentica a conexão, registra
a sessão, processa comandos em laço e limpa tudo ao final.

\begin{enumerate}
  \item Ler a primeira linha e conferir a credencial. Se falhar, responder
        \texttt{LOGIN\_FAIL} e fechar a conexão.
  \item Registrar a sessão na tabela de sessões, para o \textit{broadcast}
        encontrá-la.
  \item Em laço: ler uma linha, identificar o comando pela primeira palavra e
        chamar o tratador correspondente.
  \item Ao receber \texttt{SAIR}, ou quando o \texttt{recv()} indicar que a
        conexão caiu, remover a sessão e fechar o socket.
\end{enumerate}

\subsubsection*{Inserção e \textit{broadcast}}

O tratador de \texttt{ADD} valida os dados, insere na fila e avisa os demais. A
ordem importa: a inserção acontece antes da notificação, para nenhum cliente ser
avisado de um registro que ainda não existe.

\begin{enumerate}
  \item Separar sequência, identificador e nome da linha recebida.
  \item Validar: identificador positivo, nome não vazio e dentro do limite de
        49 caracteres.
  \item Conferir o número de sequência. Se for repetido, devolver a resposta
        guardada e encerrar aqui.
  \item Inserir na fila, com o \textit{mutex} travado.
  \item Gravar no histórico e reescrever o retrato da fila.
  \item Enviar o \textit{broadcast} a todas as sessões, exceto à de origem.
  \item Responder \texttt{ADD\_OK} ao autor e guardar essa resposta.
\end{enumerate}

\subsubsection*{\textit{Heartbeat}}

O algoritmo compara o marcador da conexão com o tamanho atual da fila. Se o
tamanho for maior, o servidor devolve os registros do intervalo e avança o
marcador. Se não, devolve \texttt{ALIVE}.

Há um detalhe necessário para a regra funcionar. Quando o próprio cliente faz um
\texttt{ADD}, seu marcador avança na hora. Sem isso, o \textit{heartbeat}
devolveria a ele o usuário que ele mesmo acabou de cadastrar, o que contraria a
definição de ``usuários cadastrados por outros clientes''.

\subsection{Decisões de implementação}

{\footnotesize
\begin{longtable}{L{4.6cm} L{8.6cm}}
\caption{Principais decisões tomadas durante a codificação.}\\
\toprule
\textbf{Decisão} & \textbf{Motivo} \\
\midrule
\endfirsthead
\toprule
\textbf{Decisão} & \textbf{Motivo} \\
\midrule
\endhead
O cliente também usa duas threads: uma lê o socket, a outra lê o teclado &
Com um fluxo só, o programa ficaria preso no \texttt{fgets} e a notificação de
\textit{broadcast} só apareceria quando o operador digitasse algo. Isso
contraria a exigência de atualização em tempo real. A thread receptora é a única
que chama \texttt{recv()} \\
\addlinespace
O texto que trafega na rede é ASCII puro; a acentuação existe só na exibição &
O enunciado especifica texto ASCII e o tráfego será comparado com esta
documentação. Bytes iguais em qualquer máquina evitam divergência por
codificação. O cliente reescreve a mensagem com acento apenas ao mostrá-la na
tela \\
\addlinespace
A fila é lida em lotes de 32 registros ao ser enviada &
Mantém o \textit{mutex} travado por intervalos curtos. Uma listagem longa não
impede outros clientes de inserir \\
\addlinespace
O nome é validado uma única vez, no servidor &
Corrige um defeito encontrado em teste. O nome era cortado ao entrar na fila,
mas a confirmação e o histórico usavam o texto original. Hoje fila, histórico,
\textit{broadcast} e resposta usam exatamente o mesmo texto \\
\addlinespace
O socket vai para a thread dentro de uma \texttt{struct} alocada, não como
inteiro convertido em ponteiro &
Converter \texttt{int} em \texttt{void*} não é portável em 64 bits. A
\texttt{struct} ainda permite levar o IP de origem e o número da conexão junto \\
\addlinespace
As threads são criadas já desatachadas &
Equivale a chamar \texttt{pthread\_detach} logo após a criação, com uma chamada a
menos por conexão. Faz diferença em 10.000 conexões \\
\addlinespace
O servidor não tem encerramento ordenado &
Ele roda até ser interrompido. Como toda gravação é seguida de \texttt{fflush},
nada se perde, e o sistema operacional libera os recursos do processo. Evita
código que nunca é executado \\
\bottomrule
\end{longtable}}

O código passou por uma revisão deliberada depois de funcionar. Sete mecanismos
que existiam por precaução, e não por necessidade, foram removidos, entre eles o
\texttt{pthread\_rwlock\_t} da tabela de sessões e o tratador de
\texttt{SIGINT}. O programa saiu de 1.545 para 1.326 linhas de código sem perder
nenhuma funcionalidade, e toda a bateria de testes foi repetida com resultados
idênticos.

% =============================================================================
% ITEM 3 - DECISOES OMISSAS NA ESPECIFICACAO
% =============================================================================
\section{Decisões de implementação omissas na especificação}
\label{sec:omissos}

O enunciado exige que qualquer elemento não especificado, mas necessário para o
protocolo funcionar, seja descrito de forma explícita. Esta seção é essa
declaração. Ela cobre quatro grupos: o que foi acrescentado, o que foi decidido
diante de ambiguidade, o que ficou fora de escopo e as limitações conhecidas.

\subsection{Elementos acrescentados por necessidade do protocolo}

\begin{enumerate}
  \item \textbf{Tabela de sessões} (\texttt{sessoes.c}). O \textit{broadcast}
        exige saber, a cada instante, quais sockets estão conectados e
        autenticados. Implementada como vetor de tamanho fixo
        (\texttt{MAX\_SESSOES = 12000}), protegido por um \textit{mutex}.

  \item \textbf{Serialização de escrita por socket.} Sem ela, uma mensagem de
        \textit{broadcast} disparada por outra thread poderia entrar no meio de
        uma resposta de várias linhas e quebrar o protocolo. Cada sessão tem um
        \textit{mutex} de envio próprio, mantido travado do cabeçalho ao rodapé
        do bloco.

  \item \textbf{Regra para evitar impasse (\textit{deadlock}).} O \textit{mutex}
        de envio de uma sessão é obtido sem travar a tabela. Quem o obtém é a
        própria thread dona da sessão, e só ela remove aquela entrada, então a
        entrada existe com certeza. A regra que sustenta isso: uma thread nunca
        remove a própria sessão enquanto segura o \textit{mutex} de envio dela.

  \item \textbf{Sinal \texttt{SIGPIPE} ignorado}, no cliente e no servidor.
        Escrever em um socket que o outro lado já fechou gera esse sinal, cujo
        comportamento padrão é encerrar o processo inteiro. Ignorá-lo transforma
        a situação em um erro de retorno, que o código já trata.

  \item \textbf{Pilha reduzida das threads}, de 8~MB para 256~KB, definida por
        \texttt{pthread\_attr\_setstacksize}. Com o valor padrão, 10.000 threads
        reservariam 80~GB de espaço de endereçamento e o teste de carga seria
        impossível.

  \item \textbf{Numeração das conexões nas linhas de log.} As mensagens começam
        exatamente com o texto das telas do enunciado e recebem, no fim, o
        número sequencial e a origem --- por exemplo,
        \texttt{Novo cliente conectado. [\#4271] 192.168.0.20:51344}. Sem isso,
        o print do teste de 10.000 clientes seria uma parede de linhas
        idênticas, e o enunciado exige que seja possível identificar todas as
        conexões.
\end{enumerate}

\subsection{Pontos ambíguos e como foram resolvidos}

{\footnotesize
\begin{longtable}{L{4.6cm} L{8.6cm}}
\caption{Pontos em aberto no enunciado e a decisão adotada.}\\
\toprule
\textbf{Ponto} & \textbf{Decisão} \\
\midrule
\endfirsthead
\toprule
\textbf{Ponto} & \textbf{Decisão} \\
\midrule
\endhead
\textit{Heartbeat}: texto da Observação contra o \textit{print} &
A Observação diz que a função devolve os usuários cadastrados por outros
clientes, ou \texttt{ALIVE}. O \textit{print} mostra a fila inteira. Seguimos o
texto, por ser a especificação escrita \\
\addlinespace
``Atualizar filas em tempo real'' no cliente &
Implementado como notificação assíncrona do novo usuário, não como réplica local
da fila. É o que os \textit{prints} mostram. O operador pede a fila completa com
a opção \texttt{2} quando quiser \\
\addlinespace
Terminologia ``paciente'' contra ``usuário'' &
Adotado ``usuário'', termo usado no texto do enunciado \\
\addlinespace
Persistência: banco de dados ou arquivo &
Arquivo de texto, opção permitida pelo enunciado. Quatro arquivos, descritos na
Seção~\ref{sec:persistencia} \\
\addlinespace
Endereço do servidor &
Constante \texttt{SERVER\_IP} em \texttt{comum.h}, alterada e recompilada ao
mudar de máquina. Nenhum código visto na disciplina faz descoberta dinâmica de
rede \\
\addlinespace
Autenticação &
Usuário e senha fixos no código, sem validação contra arquivo ou banco. Uma
credencial única, coerente com o sistema ter um único tipo de operador \\
\addlinespace
Unicidade de identificador &
Não há verificação. Dois usuários podem ter o mesmo identificador, e nada no
enunciado exige o contrário \\
\addlinespace
Capacidade da fila &
Limite fixo de 20.000 registros. Ao encher, o servidor responde
\texttt{ERRO Fila cheia} sem derrubar nada \\
\addlinespace
Tamanho do nome &
Até 49 caracteres. Acima disso o cadastro é recusado, não cortado \\
\addlinespace
Identificador inválido &
Recusado se for menor ou igual a zero \\
\addlinespace
Recuperação da fila &
O servidor sempre inicia com a fila vazia e não recarrega \texttt{fila.txt}. Os
arquivos de log e histórico continuam existindo entre execuções \\
\bottomrule
\end{longtable}}

A recusa do nome longo merece uma justificativa. Cortar o nome traria dois
problemas: o cliente receberia a confirmação de um nome diferente do que enviou,
e o corte em 49 bytes poderia cair no meio de um caractere acentuado, que ocupa
dois bytes. Nomes com acento são aceitos normalmente.

\subsection{Escopo declarado: o que não foi implementado}

A introdução do enunciado cita funcionalidades que não aparecem nas telas, na
Observação nem nos oito itens exigidos desta documentação. Elas ficaram fora de
escopo, e a decisão está declarada aqui, como o enunciado exige.

{\footnotesize
\begin{longtable}{L{4.6cm} L{8.6cm}}
\caption{Funcionalidades citadas apenas na introdução e não implementadas.}\\
\toprule
\textbf{Item} & \textbf{Motivo} \\
\midrule
\endfirsthead
\toprule
\textbf{Item} & \textbf{Motivo} \\
\midrule
\endhead
Painéis administrativos &
Citado apenas na introdução. Exigiria um tipo de cliente com interface própria,
que não aparece em nenhuma tela \\
\addlinespace
Balanceamento de carga &
Pressupõe mais de um servidor. A arquitetura do enunciado mostra um servidor
central único \\
\addlinespace
Métricas em tempo real para administradores &
Depende do painel administrativo \\
\addlinespace
Papéis distintos de recepcionista, coordenador e administrador &
O sistema tem um único tipo de operador, coerente com a credencial única
prevista para o login \\
\addlinespace
Posição numérica do usuário na fila &
A resposta de \texttt{LIST} segue o formato exato do \textit{print}
(\texttt{id - nome}), sem número de posição. A ordem das linhas já é a ordem de
atendimento \\
\addlinespace
Acompanhamento via aplicativo ou página web &
O enunciado determina cliente e servidor em C, iniciados por linha de comando \\
\bottomrule
\end{longtable}}

\subsection{Limitação conhecida: estouro de inteiro no identificador}

O identificador é lido com \texttt{sscanf(entrada, "\%d", \&id)}. Um número acima
do limite do tipo \texttt{int} sofre estouro silencioso e é gravado com outro
valor. Em teste, \texttt{99999999999999} entrou na fila como
\texttt{276447231}.

Identificadores negativos e zero são recusados, mas o estouro produz um valor
positivo, então passa pela validação. Corrigir exigiria voltar a
\texttt{strtol} com verificação de \texttt{errno}. Como é um comportamento comum
em C e sem efeito sobre o funcionamento do sistema, a limitação foi registrada
em vez de contornada.

\subsection{Persistência em arquivos de texto}
\label{sec:persistencia}

O enunciado permite escolher entre banco de dados e arquivo. A escolha foi
arquivo de texto, e os dados foram separados em quatro, distinguindo o que é do
\textbf{cliente} do que é do \textbf{usuário}.

\begin{table}[H]
\centering
\caption{Arquivos gerados pelo servidor, no subdiretório \texttt{dados/}.}
\footnotesize
\begin{tabular}{L{3.4cm} L{5.6cm} L{4.2cm}}
\toprule
\textbf{Arquivo} & \textbf{Conteúdo} & \textbf{Quando grava} \\
\midrule
\texttt{sessoes.log} & \texttt{LOGIN <ip> <data e hora>} e \texttt{LOGOUT <ip> <data e hora>} & A cada entrada e saída de cliente \\
\texttt{servidor.log} & Os mesmos eventos que aparecem na tela do servidor & A cada evento de conexão \\
\texttt{fila.txt} & Retrato atual da fila completa & Reescrito a cada \texttt{ADD} \\
\texttt{historico.txt} & Registro permanente: \texttt{<data e hora> ADD <id> <nome>} & A cada \texttt{ADD} \\
\bottomrule
\end{tabular}
\end{table}

O arquivo \texttt{historico.txt} atende a dois dos itens pedidos pelo enunciado
ao mesmo tempo. Cada linha registra um usuário cadastrado com data e hora,
servindo tanto como cadastro de usuários quanto como histórico.

% =============================================================================
% ITEM 4 - RETRANSMISSAO
% =============================================================================
\section{Retransmissão de mensagens}
\label{sec:retransmissao}

\subsection{O que precisava ser garantido}

O TCP garante entrega ordenada e sem perda enquanto a conexão está de pé. Ele
garante que os bytes chegaram à outra máquina, mas não que a aplicação do outro
lado processou o pedido. Se o servidor estiver vivo e travado, ou se a conexão
cair no meio de uma operação, o TCP não resolve o problema do cliente.

A retransmissão implementada é, portanto, do nível da aplicação. Ela responde a
duas perguntas: como o cliente sabe que seu comando foi processado, e o que ele
faz quando não sabe.

\subsection{Mecanismo adotado}

O mecanismo tem três partes, e nenhuma delas usa uma mensagem de confirmação
separada.

\vspace{6pt}
\noindent\textbf{1. A resposta é a confirmação.} Toda resposta listada na
Seção~\ref{sec:protocolo} funciona como confirmação do pedido correspondente. O
\texttt{ADD\_OK} confirma que o \texttt{ADD} foi recebido e processado. Não
existe uma mensagem \texttt{ACK} à parte.

\vspace{6pt}
\noindent\textbf{2. Tempo limite e reenvio.} Depois de enviar um comando, o
cliente espera a resposta por até \textbf{3 segundos}. Se o prazo estourar, ele
reenvia o mesmo comando, até um total de \textbf{3 tentativas}. Se as três
falharem, avisa o operador com \texttt{Servidor não respondeu, tente novamente.}
e volta ao menu.

\vspace{6pt}
\noindent\textbf{3. Número de sequência e idempotência.} Cada comando carrega um
número de sequência próprio daquela conexão. O servidor guarda o último número
processado. Se receber um número repetido, ele devolve a mesma resposta de antes
\textbf{sem processar de novo}. É isso que impede um reenvio de cadastrar o
mesmo usuário duas vezes.

Do lado do cliente há uma regra complementar. Uma resposta só é exibida se
houver um comando esperando por ela. Cópias atrasadas, que chegam depois de o
cliente ter desistido, são descartadas em silêncio. Sem essa regra, uma cópia
atrasada seria confundida com a resposta do comando seguinte.

\subsection{Avaliação honesta do mecanismo}

Sobre uma única conexão TCP, o reenvio é em boa parte redundante, porque o
próprio TCP já retransmite segmentos perdidos. Vale reconhecer isso.

O que de fato agrega é o \textbf{número de sequência}. Ele é o que torna a
operação idempotente e impede que um comando repetido produza efeito duplicado.
Essa garantia o TCP não oferece, porque ela é sobre o processamento na
aplicação, não sobre a entrega dos bytes.

% =============================================================================
% ITEM 5 - TESTES E ANALISE
% =============================================================================
\section{Testes e análise}

\subsection{Como os testes foram feitos}

Os testes seguiram três frentes, todas repetidas a cada alteração do código:

\begin{enumerate}
  \item \textbf{Teste funcional} --- dois clientes reproduzindo a sequência das
        telas do enunciado.
  \item \textbf{Casos especiais} --- um programa em Python conversando direto com
        o servidor pelo socket, enviando comandos válidos e inválidos e
        conferindo cada resposta. Falar o protocolo em texto permite testar
        situações que o \texttt{./cliente} não deixaria acontecer, como um
        comando malformado ou um comando antes do login.
  \item \textbf{Carga} --- o programa \texttt{carga.c}, com 100, 1.000 e 10.000
        clientes.
\end{enumerate}

\subsection{Teste funcional}

\begin{table}[H]
\centering
\caption{Teste funcional com dois clientes conectados ao mesmo servidor.}
\footnotesize
\begin{tabular}{c L{4.4cm} L{6.0cm} c}
\toprule
\textbf{\#} & \textbf{Ação} & \textbf{Resultado esperado} & \textbf{Obtido} \\
\midrule
1 & Clientes A e B conectam & \texttt{Conectado ao servidor.} e \texttt{Servidor: LOGIN\_OK} & OK \\
2 & A cadastra \texttt{10/Abel} e \texttt{30/Breno} & \texttt{Usuário Abel adicionado.} & OK \\
3 & B recebe as notificações & \texttt{[Broadcast] Novo usuário: Abel} e \texttt{Breno} & OK \\
4 & B cadastra \texttt{20/Carlos} e \texttt{40/Dario} & Confirmação em B, \textit{broadcast} em A & OK \\
5 & A escolhe \texttt{2} & Os quatro usuários, entre cabeçalho e rodapé & OK \\
6 & A escolhe \texttt{3} & Só \texttt{20 - Carlos} e \texttt{40 - Dario}, cadastrados por B & OK \\
7 & B escolhe \texttt{3} & \texttt{Heartbeat: ALIVE}, pois B só cadastrou os próprios & OK \\
8 & Ambos escolhem \texttt{0} & \texttt{Cliente desconectado.} na tela do servidor & OK \\
\bottomrule
\end{tabular}
\end{table}

O passo 7 confere com o \textit{print} do enunciado, que também mostra
\texttt{ALIVE} para o cliente que acabou de cadastrar. O passo 6 é onde seguimos
o texto da Observação em vez do \textit{print}: o \textit{print} mostra a fila
inteira, e o sistema devolve apenas os usuários de outros clientes.

\subsection{Casos especiais: 30 de 30 aprovados}

Foram testados 30 casos especiais, agrupados por tema. Todos passaram. Em todos
os casos de comando inválido, \textbf{a conexão permanece aberta} --- um erro do
operador não derruba o cliente.

{\footnotesize
\begin{longtable}{c L{6.0cm} L{6.0cm}}
\caption{Casos especiais testados e resposta obtida.}\\
\toprule
\textbf{\#} & \textbf{Situação testada} & \textbf{Resposta obtida} \\
\midrule
\endfirsthead
\toprule
\textbf{\#} & \textbf{Situação testada} & \textbf{Resposta obtida} \\
\midrule
\endhead
\multicolumn{3}{l}{\textit{Autenticação}}\\
\addlinespace
1 & Senha errada & \texttt{LOGIN\_FAIL} \\
2 & Conexão após falha de login & Encerrada pelo servidor \\
3 & \texttt{LOGIN} sem a senha & \texttt{LOGIN\_FAIL} \\
4 & Comando antes do login & \texttt{LOGIN\_FAIL} \\
5 & Credencial correta & \texttt{LOGIN\_OK} \\
\addlinespace
\multicolumn{3}{l}{\textit{Comandos malformados}}\\
\addlinespace
6 & Comando inexistente & \texttt{ERRO Comando desconhecido} \\
7 & \texttt{ADD} sem argumento & \texttt{ERRO Formato esperado...} \\
8 & \texttt{ADD} só com a sequência & \texttt{ERRO Formato esperado...} \\
9 & \texttt{ADD} com identificador não numérico & \texttt{ERRO Formato esperado...} \\
10 & \texttt{ADD} sem o nome & \texttt{ERRO Nome invalido...} \\
11 & \texttt{LIST} sem número de sequência & \texttt{ERRO Formato esperado...} \\
12 & \texttt{HEARTBEAT} sem número de sequência & \texttt{ERRO Formato esperado...} \\
13 & Linha em branco & Ignorada; a conexão segue \\
\addlinespace
\multicolumn{3}{l}{\textit{Validação de dados}}\\
\addlinespace
14 & Identificador negativo & \texttt{ERRO O identificador deve ser um numero positivo} \\
15 & Identificador zero & \texttt{ERRO ...numero positivo} \\
16 & Nome com 60 caracteres & \texttt{ERRO Nome invalido, vazio ou longo demais} \\
17 & Nome com espaços & Aceito \\
18 & Nome com acentos & Aceito \\
\addlinespace
\multicolumn{3}{l}{\textit{Retransmissão, lado do servidor}}\\
\addlinespace
19 & \texttt{ADD} repetido com o mesmo número de sequência & Mesma resposta \texttt{ADD\_OK 10 Abel} \\
20 & Fila após o \texttt{ADD} repetido & Um único registro; não duplicou \\
21 & Sequência menor que a última processada & \texttt{ERRO Numero de sequencia invalido} \\
22 & Sequência zero & \texttt{ERRO Numero de sequencia invalido} \\
\addlinespace
\multicolumn{3}{l}{\textit{Heartbeat e broadcast}}\\
\addlinespace
23 & \textit{Broadcast} chega ao outro cliente & \texttt{[Broadcast] Novo usuario: Breno} \\
24 & \textit{Heartbeat} com novidade & Só \texttt{30 - Breno}, cadastrado por outro \\
25 & O mesmo \textit{heartbeat} reenviado & Bloco idêntico ao anterior \\
26 & \textit{Heartbeat} seguinte, sem novidade & \texttt{ALIVE} \\
27 & Esse \textit{heartbeat} reenviado & \texttt{ALIVE} de novo \\
28 & Novo cadastro e novo \textit{heartbeat} & Só \texttt{40 - Dario}, o registro novo \\
29 & Esse \textit{heartbeat} reenviado & Mesmo bloco \\
\addlinespace
\multicolumn{3}{l}{\textit{Falha de conexão}}\\
\addlinespace
30 & Cliente cai sem enviar \texttt{SAIR} & O servidor detecta, remove a sessão e continua atendendo os demais \\
\bottomrule
\end{longtable}}

\vspace{6pt}
\noindent\textbf{Análise.} Os casos 25, 27 e 29 são os mais importantes do
conjunto. Eles comprovam a idempotência: reenviar um comando devolve exatamente
a mesma resposta e não avança o marcador do que já foi visto. Sem essa
propriedade, o mecanismo de retransmissão da Seção~\ref{sec:retransmissao} seria
inseguro, porque cada reenvio produziria um efeito novo.

O caso 30 comprova que a queda abrupta de um cliente é tratada como um evento
normal. O \texttt{recv()} devolve zero quando o outro lado fecha, e a thread sai
do laço, remove a sessão e fecha o socket, sem afetar as demais.

\subsection{Retransmissão vista do lado do cliente}

O comportamento do cliente foi testado contra um servidor de mentira, que aceita
o login e depois não responde a mais nada. O objetivo era forçar o tempo limite.

\begin{figure}[H]
\centering
\begin{minipage}{0.75\textwidth}
\ttfamily\footnotesize
\begin{verbatim}
t = 0,0s   servidor recebeu: ADD 1 10 Abel
t = 3,0s   servidor recebeu: ADD 1 10 Abel
t = 6,0s   servidor recebeu: ADD 1 10 Abel
\end{verbatim}
\end{minipage}
\caption{Registro do servidor de mentira durante o teste de retransmissão.}
\end{figure}

Os resultados confirmaram o projetado:

\begin{itemize}
  \item exatamente 3 envios, um a cada 3 segundos;
  \item os três com o mesmo número de sequência, que é o que permite ao servidor
        reconhecer o reenvio;
  \item o cliente avisou o operador ao fim das tentativas;
  \item o servidor de mentira enviou depois 3 cópias atrasadas da resposta, e o
        cliente ignorou todas, porque já não havia comando esperando;
  \item o comando seguinte funcionou normalmente, sem confusão.
\end{itemize}

\subsection{Compilação}

\begin{table}[H]
\centering
\caption{Verificações de compilação.}
\small
\begin{tabular}{L{7.6cm} L{5.4cm}}
\toprule
\textbf{Verificação} & \textbf{Resultado} \\
\midrule
\texttt{make} sem parâmetros & Gera exatamente \texttt{cliente} e \texttt{servidor} \\
Avisos com \texttt{-Wall -Wextra} & Nenhum \\
Avisos com \texttt{-Wpedantic -Wshadow -Wconversion -Wsign-conversion -Wcast-align -Wnull-dereference} & Nenhum \\
Compilação com \texttt{-std=c99}, \texttt{-std=c11} e \texttt{-std=c17}, com e sem \texttt{-O2} & Sem erro e sem aviso nas oito combinações \\
\bottomrule
\end{tabular}
\end{table}

% =============================================================================
% ITEM 6 - CARGA
% =============================================================================
\section{Teste de carga com 100, 1.000 e 10.000 clientes}

\subsection{Como a carga foi gerada}

A geração é automática, feita pelo programa \texttt{carga.c} --- o segundo
código cliente que o enunciado permite expressamente. Ele repete
\texttt{socket()}, \texttt{connect()} e login $N$ vezes, e recebe a quantidade
como parâmetro:

\begin{lstlisting}[caption={Comandos usados no teste de carga.},language=bash]
make carga
./carga 100
./carga 1000
./carga 10000
\end{lstlisting}

Duas providências foram necessárias para o teste medir o que devia medir. A
primeira: as conexões ficam \textbf{todas abertas ao mesmo tempo} até o operador
pressionar ENTER. Sem isso, o programa estaria abrindo e fechando conexões em
sequência, e o servidor nunca precisaria sustentar mais de uma por vez. A
segunda: cada cliente gerado \textbf{também se autentica}, ocupando uma sessão
no servidor exatamente como um \texttt{./cliente}.

O servidor foi reiniciado antes de cada teste, para a numeração das conexões
começar em \texttt{[\#1]}. Antes do teste de 10.000, o limite de descritores de
arquivo do terminal foi elevado com \texttt{ulimit -n 20000}.

\subsection{Resultados}

\begin{table}[H]
\centering
\caption{Resultado dos três testes de carga.}
\footnotesize
\begin{tabular}{r r r l r}
\toprule
\textbf{Clientes} & \textbf{Conexões abertas} & \textbf{Falhas} & \textbf{Numeração no servidor} & \textbf{Sessões autenticadas} \\
\midrule
100 & 100 & 0 & \texttt{[\#1]} a \texttt{[\#100]} & 100 \\
1.000 & 1.000 & 0 & \texttt{[\#1]} a \texttt{[\#1000]} & 1.000 \\
10.000 & 10.000 & 0 & \texttt{[\#1]} a \texttt{[\#10000]} & 10.000 \\
\bottomrule
\end{tabular}
\end{table}

Depois de cada teste, um cliente comum foi executado e operou normalmente. Isso
mostra que o servidor não ficou em estado ruim após sustentar a carga.

\tela{carga-100.png}{Tela do servidor durante o teste com 100 clientes, com as conexões numeradas de \texttt{[\#1]} a \texttt{[\#100]}.}{0.9}

\tela{carga-1000.png}{Tela do servidor durante o teste com 1.000 clientes.}{0.9}

\tela{carga-10000.png}{Tela do servidor durante o teste com 10.000 clientes.}{0.9}

\subsection{Análise dos resultados}

\noindent\textbf{Por que 10.000 threads couberam.} Cada thread é criada com
pilha de 256~KB, definida por \texttt{pthread\_attr\_setstacksize}. Com o padrão
do sistema, de 8~MB, 10.000 threads reservariam 80~GB de espaço de
endereçamento e o teste seria impossível. Com 256~KB, o total fica em torno de
2,5~GB de memória virtual, que é espaço reservado e não ocupado fisicamente.

\vspace{6pt}
\noindent\textbf{Por que as linhas não saem em ordem numérica perfeita.} No
teste de 100 clientes, a linha \texttt{[\#99]} apareceu depois da
\texttt{[\#100]}. Isso não é defeito. Cada conexão é atendida por uma thread
diferente, e a ordem em que elas escrevem no terminal depende de qual thread o
sistema escalona primeiro. O número é atribuído no momento da aceitação, então a
contagem está correta. Na prática, é uma evidência visual de que o servidor é
mesmo concorrente.

% =============================================================================
% ITEM 7 - TELAS DE FUNCIONAMENTO
% =============================================================================
\section{Telas do funcionamento correto}

Esta seção mostra o sistema em operação. As capturas seguem a mesma sequência do
teste funcional descrito na seção anterior.

\tela{01-servidor-iniciado.png}{Servidor recém-iniciado, mostrando \texttt{Servidor iniciado na porta 8080} e as conexões dos dois clientes.}{0.9}

\tela{02-cliente-a-login-menu.png}{Cliente A conectado e autenticado, com o menu de operação.}{0.85}

\tela{03-cliente-a-cadastro.png}{Cliente A cadastrando os usuários \texttt{10 - Abel} e \texttt{30 - Breno}.}{0.85}

\tela{04-cliente-b-broadcast.png}{Cliente B recebendo a notificação de \textit{broadcast} sem ter pedido nada. Esta é a evidência visual do funcionamento em tempo real.}{0.85}

\tela{05-ver-fila.png}{Opção \texttt{2}, com a fila completa entre o cabeçalho e o rodapé.}{0.85}

\tela{06-heartbeat-novidade.png}{Opção \texttt{3} no cliente A: apenas os usuários cadastrados pelo cliente B.}{0.85}

\tela{07-heartbeat-alive.png}{Opção \texttt{3} repetida, sem novidade desde a última verificação: resposta \texttt{ALIVE}.}{0.85}

\tela{08-caso-de-erro.png}{Tratamento de erro: cadastro recusado por identificador inválido. A conexão permanece aberta.}{0.85}

\tela{09-persistencia.png}{Conteúdo de \texttt{dados/fila.txt} e \texttt{dados/historico.txt}, gerados pelo servidor.}{0.9}

% =============================================================================
% ITEM 8 - CONCLUSAO E REFERENCIAS
% =============================================================================
\section{Conclusão}

O sistema atende ao que o enunciado especifica. Cliente e servidor reproduzem as
telas de exemplo ponta a ponta, o \textit{heartbeat} segue a Observação, a
persistência grava os cinco tipos de dado pedidos e o servidor sustentou 10.000
conexões simultâneas autenticadas sem nenhuma falha. O código final tem 1.326
linhas em seis módulos e compila sem nenhum aviso.

Mais interessante que o resultado, porém, é o caminho até ele. Quatro
dificuldades marcaram o desenvolvimento.

\vspace{6pt}
\noindent\textbf{A delimitação de mensagens.} Nenhum dos códigos vistos em aula
precisa resolver esse problema, porque cada um troca uma única mensagem por
conexão. Aqui a conexão é persistente e carrega dezenas de mensagens, e o TCP
não separa uma da outra. Foi preciso projetar do zero a convenção de uma linha
por mensagem e o buffer que junta os pedaços recebidos. Esse foi o ponto do
trabalho que mais exigiu leitura fora do material da disciplina.

\vspace{6pt}
\noindent\textbf{A escolha entre \texttt{fork()} e threads.} O \texttt{fork()} é
a técnica mais detalhada em aula, e a primeira escolha natural. O problema
apareceu ao pensar na fila: processos têm memória isolada, e compartilhar a fila
exigiria comunicação entre processos, que não foi ensinada. As threads resolvem
o compartilhamento sozinhas e cobram o \textit{mutex} em troca. Foi uma troca
consciente, não uma preferência.

\vspace{6pt}
\noindent\textbf{O teste de 10.000 clientes não funcionou de primeira.} A
primeira tentativa falhou por falta de memória. A causa levou algum tempo para
aparecer: cada thread reserva 8~MB de pilha por padrão, o que daria 80~GB. A
solução foi reduzir a pilha para 256~KB. Esse foi o momento em que ficou claro
que ``suportar $N$ clientes'' depende de detalhes que não estão no código do
protocolo.

\vspace{6pt}
\noindent\textbf{Dois defeitos apareceram nos próprios testes.} O primeiro: o
nome era cortado ao entrar na fila, mas a confirmação e o histórico usavam o
texto original, então o cliente recebia a confirmação de um nome diferente do
que enviou. Passou a ser recusado em vez de cortado. O segundo: as conexões do
teste de carga apareciam como falha de autenticação quando, na verdade,
encerravam antes de tentar logar. São situações diferentes e passaram a ser
registradas como tais. Nenhum dos dois teria sido encontrado sem os casos
especiais.

\vspace{6pt}
\noindent\textbf{O código foi enxugado depois de funcionar.} Sete mecanismos que
existiam por precaução, e não por necessidade, foram removidos após a primeira
versão passar em todos os testes. O programa saiu de 1.545 para 1.326 linhas sem
perder nenhuma funcionalidade, e a bateria de testes foi repetida por inteiro com
resultados idênticos. O critério foi manter apenas o que o enunciado exige ou o
que o protocolo precisa para funcionar.

\vspace{6pt}
\noindent\textbf{Uma ressalva sobre a retransmissão.} Sobre uma conexão TCP
única, o reenvio de comandos é em boa parte redundante, porque o TCP já
retransmite segmentos perdidos. O que de fato agrega valor é o número de
sequência, que impede um comando repetido de cadastrar o mesmo usuário duas
vezes. Reconhecer isso parece mais correto do que apresentar o mecanismo como
indispensável.

O trabalho deixa caminhos abertos. Ler o endereço do servidor de um arquivo de
configuração, em vez de recompilar, é a melhoria mais imediata. Substituir o
vetor fixo por uma estrutura dinâmica removeria o limite de 20.000 registros. E
migrar de uma thread por conexão para multiplexação com \texttt{select()}
reduziria o custo por cliente, ao preço de um código bem mais complexo.

% =============================================================================
% REFERENCIAS
% =============================================================================
\section*{Referências bibliográficas}
\addcontentsline{toc}{section}{Referências bibliográficas}

\begin{thebibliography}{9}
\raggedright

\bibitem{make}
FREE SOFTWARE FOUNDATION. \textit{GNU Make Manual}. Disponível em:
\url{http://www.gnu.org/software/make/manual/make.html}. Acesso em: 28 jul.
2026.

\bibitem{maketutor}
COLBY COLLEGE. \textit{Makefile Tutorial}. Disponível em:
\url{http://www.cs.colby.edu/maxwell/courses/tutorials/maketutor/}. Acesso em:
28 jul. 2026.

\bibitem{manpages}
THE LINUX MAN-PAGES PROJECT. \textit{Linux Programmer's Manual}. Páginas
consultadas: \texttt{socket(2)}, \texttt{bind(2)}, \texttt{listen(2)},
\texttt{accept(2)}, \texttt{send(2)}, \texttt{recv(2)},
\texttt{setsockopt(2)}, \texttt{pthread\_create(3)},
\texttt{pthread\_mutex\_init(3)}, \texttt{pthread\_cond\_timedwait(3)},
\texttt{pthread\_attr\_setstacksize(3)}, \texttt{sscanf(3)}. Disponível em:
\url{https://man7.org/linux/man-pages/}. Acesso em: 28 jul. 2026.

\bibitem{aula}
FERREIRA, Rubens Barbosa. \textit{Códigos de exemplo da disciplina de Redes de
Computadores}: \texttt{bind.c}, \texttt{escrevendo.c}, \texttt{fork.c},
\texttt{multithread.c}, \texttt{porta.c}, \texttt{redes.c}. Dourados:
Universidade Estadual de Mato Grosso do Sul, 2026.

% ==================================================================
% [CONFERIR - Victor]
% Duas coisas nesta lista dependem de voce:
%
% 1. O nome completo do professor na referencia [4]. Coloquei
%    "Rubens Barbosa Ferreira" como suposicao a partir do usuario
%    "rubensbf" que aparece nos prints do enunciado. CONFIRME o nome
%    correto antes de entregar - nome errado em referencia e pior que
%    referencia ausente.
%
% 2. Se voce consultou algum livro (Tanenbaum, Kurose, Stevens) ou
%    algum slide da disciplina durante o desenvolvimento, acrescente
%    aqui. Se NAO consultou, nao invente: a lista acima ja cobre as
%    fontes que o enunciado pede e as que o codigo realmente usou.
% ==================================================================

\end{thebibliography}

% =============================================================================
% MARCOS PENDENTES NESTE DOCUMENTO
%
% [A ESCREVER - Victor] Capturas de tela. Coloque os arquivos em figuras/
% com os nomes abaixo. Enquanto nao existirem, o PDF compila e mostra uma
% moldura no lugar de cada um.
%
%   figuras/01-servidor-iniciado.png
%   figuras/02-cliente-a-login-menu.png
%   figuras/03-cliente-a-cadastro.png
%   figuras/04-cliente-b-broadcast.png
%   figuras/05-ver-fila.png
%   figuras/06-heartbeat-novidade.png
%   figuras/07-heartbeat-alive.png
%   figuras/08-caso-de-erro.png
%   figuras/09-persistencia.png
%   figuras/carga-100.png
%   figuras/carga-1000.png
%   figuras/carga-10000.png
%
% [CONFERIR - Victor] Referencia [4]: nome completo do professor.
% =============================================================================

\end{document}
